\documentclass[12pt,a4paper]{article}

%% PACHETE MATEMATICE
\usepackage{amsmath,amssymb,amsfonts}
\usepackage{amsthm}
\usepackage{mathpazo}
\usepackage{eulervm}
\usepackage{enumerate}
\usepackage{color}
\usepackage{graphicx}
\usepackage[all]{xy}
\usepackage{xcolor}
\usepackage{framed}
\usepackage[unicode]{hyperref}
\setcounter{MaxMatrixCols}{20}
\usepackage{multicol}
% \usepackage[romanian]{babel}

% note de subsol oricît de lungi
\interfootnotelinepenalty=10000

%% ASPECT
\linespread{1.05}
\usepackage[T1]{fontenc}
\usepackage[utf8x]{inputenc}
\usepackage[margin=2cm]{geometry}
% \usepackage[color]{showkeys}
\usepackage{anyfontsize}
\usepackage{titlesec}
\titleformat*{\section}{\large\bfseries}

\usepackage{fancyhdr}
\pagestyle{fancy}
\rhead{\small \color{gray} Notițe de Adrian Manea}
\lhead{\small \color{gray} ETTI-AM2, 2017---2018, M. Joița \& A. Niță}


\usepackage[autostyle = false, style = german]{csquotes}
\usepackage[romanian]{babel}
\newcommand\qq{\enquote}

%% COMENZILE MELE
\renewcommand*{\proofname}{Dem.:}
\newcommand\bb{\mathbb}
\newcommand\dr{\mathrm}
\newcommand\kal{\mathcal}
\newcommand\fk{\mathfrak}
\newcommand\op{\oplus}
\newcommand\ot{\otimes}
\newcommand\ds{\displaystyle}
\newcommand\id{\indent}
\newcommand\nid{\noindent}
\newcommand\seq{\subseteq}
\newcommand\sneq{\subsetneq}
\newcommand\speq{\supseteq}
\newcommand\spneq{\supsetneq}
\newcommand\rar{\rightarrow}
\newcommand\Rar{\Rightarrow}
\newcommand\xrar{\xrightarrow}
\newcommand\lar{\leftarrow}
\newcommand\Lar{\Leftarrow}
\newcommand\xlar{\xleftarrow}
\newcommand\lrar{\leftrightarrow}
\newcommand\Lrar{\Leftrightarrow}
\newcommand\ideal{\trianglelefteq}
\newcommand\ai{\text{ a.\^{i}. }}
\newcommand\sk{(\textit{Schi\c{t}\u{a}}) }
\newcommand\rhup{\rightharpoonup}
\newcommand\rhdn{\rightharpoondown}
\newcommand\lhup{\leftharpoonup}
\newcommand\lhdn{\leftharpoondown}
\newcommand\vid{\emptyset}
\newcommand\wh{\widehat}
\newcommand\ol{\overline}
\newcommand\NN{\mathbb{N}}
\newcommand\RR{\mathbb{R}}
\newcommand\ZZ{\mathbb{Z}}
\newcommand\QQ{\mathbb{Q}}
\newcommand\CC{\mathbb{C}}

%% STILURILE MELE
\newtheoremstyle{dotless}{}{}{\itshape}{}{\bfseries}{:}{ }{}

\theoremstyle{dotless}
\newtheorem{theorem}{Teorem\u{a}}
\newtheorem{proposition}{Propozi\c{t}ie}
\newtheorem{lemma}{Lem\u{a}}
\newtheorem{corollary}{Corolar}
\newtheorem{exercise}{Exerci\c{t}iu}

\newtheoremstyle{dotlessdr}{}{}{}{}{\bfseries}{:}{ }{}
\theoremstyle{dotlessdr}
\newtheorem{example}{Exemplu}
\newtheorem{definition}{Defini\c{t}ie}
\newtheorem{remark}{Observa\c{t}ie}

\newenvironment{solution}
{\begin{proof}[Soluție:\nopunct]}
  {\end{proof}}


\begin{document}
% \definecolor{labelkey}{RGB}{222,0,0}

\begin{center}
  {\large\textbf{Seminar 10 \\
      Transformata Fourier}}
\end{center}

\vspace{2cm}

\section*{Integrala Fourier}
\indent\indent
Seriile Fourier sînt utile pentru dezvoltarea unor funcții periodice (sau convertibile
în unele periodice). Însă dacă funcțiile sînt arbitrare, se folosește o metodă care
extinde pe cea a seriilor Fourier, anume integralele Fourier.

Amintim că o funcție periodică $ f_L(x) $, de perioadă $ 2L $, poate fi dezvoltată în
serie Fourier cu formula:
\[
  f_L(x) = a_0 + \sum_{n \geq 1} (a_n \cos w_n x + b_n \sin w_n x), \quad %
  w_n = \frac{n\pi}{L}.
\]

În această serie punem $ L \to \infty $ și, după impunerea unor condiții
de convergență, ajungem la \emph{integrala Fourier} a funcției $ f $, anume:
\[
  f(x) = \int_0^\infty A(w) \cos wx + B(w) \sin wx dw,
\]
unde funcțiile $ A $ și $ B $ sînt date de:
\[
  A(w) = \frac{1}{\pi} \int_{-\infty}^\infty f(v) \cos wv dv, \quad %
  B(w) = \frac{1}{\pi} \int_{-\infty}^\infty f(v) \sin wv dv.
\]

Ca în cazul seriilor Fourier, dacă $ f $ este o funcție pară, atunci $ B(w) = 0 $,
iar integrala Fourier devine o integrală de cosinusuri:
\[
  f(x) = \int_0^\infty A(w) \cos wx dw, \text{ unde } %
  A(w) = \frac{2}{\pi} \int_0^\infty f(v) \cos wv dv.
\]

Similar, pentru $ f $ funcție impară, avem $ A(w) = 0 $, iar integrala Fourier devine
o integrală de sinusuri:
\[
  f(x) = \int_0^\infty B(w) \sin wx dw, \text{ unde } %
  B(w) = \frac{2}{\pi} \int_0^\infty f(v) \sin wv dv.
\]

%%%%%%%%%%%%%%%%%%%%%%%%%%%%%%%%%%%%%%%%%%%%%%%%%%%%%%%%%%%%%%%%%%%%%% 

\section*{Transformata Fourier}

\indent\indent
Ideea de bază a unei transformări Fourier este următoarea. Dacă se dă o funcție periodică
(sau convertită la una periodică printr-un artificiu de repetiție, practic), ei i se
asociază \emph{seria Fourier}. Aceasta o aproximează cu o serie de sinusuri și cosinusuri,
funcțiile periodice cel mai des întîlnite. Practic, are loc o superpoziție de termeni cu
sinusuri și cosinusuri, un fel de interferență a undelor electromagnetice. La pasul
următor, \emph{transformata Fourier} preia minimele și maximele acestor \qq{interferențe},
iar rezultatul este un semnal (aproape) discret (\qq{digitalizat}), care reprezintă
valorile cele mai importante din semnal.\footnote{O explicație animată este dată \href{https://en.wikipedia.org/wiki/Fourier_transform\#/media/File:Fourier_transform_time_and_frequency_domains_(small).gif}{aici}.}

Rescriem formula integralei Fourier, înlocuind funcțiile $ A $ și $ B $:
\[
  f(x) = \frac{1}{\pi} \int_0^\infty \int_{-\infty}^\infty f(v) \cdot %
  \big( \cos wv \cos wx + \sin wv \sin wx \big) dvdw.
\]

Acum putem folosi formule trigonometrice uzuale și observăm că putem rescrie:
\[
  f(x) = \frac{1}{\pi} \int_0^\infty \Big( \int_{-\infty}^\infty %
  f(v) \cos (wx - wv) dv \Big) dw.
\]

Cum funcția $ \cos $ este pară, integrala de la $ 0 $ la $ \infty $ este jumătate
din integrala pe tot $ \RR $, deci putem rescrie în forma:
\[
  f(x) = \frac{1}{2\pi} \int_{-\infty}^\infty \Big( \int_{-\infty}^\infty f(v) \cdot %
  \cos (wx - wv) dv \Big) dw.
\]

\begin{remark}\label{rk:fourier-sin}
  Să remarcăm că, dacă în integrala de mai sus aveam funcția $ \sin $ în loc de $ \cos $,
  integrala ar fi fost nulă, din imparitatea funcției $ \sin $.
\end{remark}

Conform observației de mai sus, putem adăuga un termen similar cu $ \sin $, fără a
schimba valoarea integralei. Acest lucru este util pentru a trece pe domeniu
complex, unde putem porni de la formula lui Euler de scriere polară a unui număr
complex. Așadar, avem:
\[
  f(v) \cos (wx - wv) + if(v) \sin (wx - wv) = f(v) e^{i(wx - wv)}.
\]

Trecînd la integrală, obținem \emph{integrala Fourier complexă:}
\begin{equation}
  \label{eq:int-four-compx}
  f(x) = \frac{1}{2\pi} \int_{-\infty}^\infty \int_{-\infty}^\infty %
  f(v) e^{iw(x - v)} dvdw.
\end{equation}

Dacă descompunem funcția exponențială într-un produs și scriem ca produs de integrale,
obținem:
\[
  f(x) = \frac{1}{\sqrt{2\pi}} \int_{-\infty}^\infty \Big[ \frac{1}{\sqrt{2\pi}} %
  \int_{-\infty}^\infty f(v) e^{-iwv} dv \Big] e^{iwx} dw.
\]

Integrala din interior este exact \textbf{transformata Fourier} a lui $ f $:
\begin{equation}
  \label{eq:transf-f}
  \widehat{f}(w) = \frac{1}{\sqrt{2\pi}} \int_{-\infty}^\infty f(x) e^{-iwx} dx.
\end{equation}

Atunci, dacă înlocuim în formula de mai sus, obținem:
\begin{equation}
  \label{eq:transf-f-inv}
  f(x) = \frac{1}{\sqrt{2\pi}} \int_{-\infty}^\infty \widehat{f}(w) e^{iwx} dw,
\end{equation}
care este formula \emph{transformării Fourier inverse}.

O altă notație pentru transformata Fourier este $ \widehat{f} = \kal{F}(f) $ și
pentru transformata inversă, $ f = \kal{F}^{-1}(\widehat{f}) $.

Pentru orice funcție care satisface anumite proprietăți, transformata Fourier există:

\begin{theorem}\label{thm:ex-four}
  Dacă $ f $ este absolut integrabilă (adică integrala funcției $ |f(x)| $ este convergentă)
  și continuă pe orice interval finit, atunci transformata Fourier dată de
  formula~\eqref{eq:transf-f} există.
\end{theorem}

Să vedem cîteva exemple.

\begin{example}
  Găsiți transformata Fourier a funcției:
  \[
    f(x) =
    \begin{cases}
      1, & |x| < 1 \\
      0, & \text{în rest}
    \end{cases}.
  \]
\end{example}

\emph{Soluție:} Folosind definiția, integrăm:
\[
  \widehat{f}(w) = \frac{1}{\sqrt{2\pi}} \int_{-1}^1 e^{-iwx} dx = %
  \frac{1}{\sqrt{2\pi}} \cdot \frac{e^{-iwx}}{-iw} \Big|_{-1}^1 = %
  \frac{1}{-iw\sqrt{2\pi}}(e^{-iw} - e^{iw}).
\]

Folosind formula lui Euler pentru $ e^{\pm iw} $, avem, în fine:
\[
  \widehat{f}(w) = \sqrt{\frac{\pi}{2}} \frac{\sin w}{w}.
\]

Un alt exemplu:

\begin{example}
  Găsiți transformata Fourier a funcției:
  \[
    f(x) =
    \begin{cases}
      e^{-ax}, & x > 0 \\
      0, & x < 0
    \end{cases}, a > 0.
  \]
\end{example}

\emph{Soluție:} Din definiție, avem:
\begin{align*}
  \widehat{f}(w) &= \frac{1}{\sqrt{2\pi}} \int_0^\infty e^{-ax} e^{-iwx} dx \\
                 &= \frac{1}{\sqrt{2\pi}} \frac{e^{-(a + iw)x}}{-(a + iw)} \Big|_{x = 0}^\infty \\
                 &= \frac{1}{\sqrt{2\pi}(a + iw)}
\end{align*}

%%%%%%%%%%%%%%%%%%%%%%%%%%%%%%%%%%%%%%%%%%%%%%%%%%%%%%%%%%%%%%%%%%%%%% 

\section*{Proprietăți ale transformatei Fourier}


\indent\indent \textbf{Liniaritate:} Transformata Fourier este o operație liniară:
\[
  \kal{F}(af + bg) = a \kal{F}(f) + b \kal{F}(g),
\]
pentru orice funcții $ f, g $ care admit transformată Fourier și $ a, b \in \RR $.

\vspace{1cm}

\textbf{Transformata derivatei:} Dacă $ f $ este o funcție continuă și $ f(x) \to 0 $
pentru $ |x| \to \infty $, iar $ f'(x) $ este absolut integrabilă, atunci:
\[
  \kal{F}(f'(x)) = iw \kal{F}(f).
\]

Mai departe, pentru derivate superioare, obținem, de exemplu:
\[
  \kal{F}(f^{\prime\prime}(x)) = -w^2 \kal{F}(f(x)).
\]

\vspace{1cm}

\textbf{Convoluție:} Fie $ f, g $ două funcții. Se definește \emph{produsul lor de %
  convoluție} $ f \ast g $ ca fiind funcția:
\[
  h(x) = (f \ast g)(x) = \int_{-\infty}^\infty f(p) g(x-p) dp = %
  \int_{-\infty}^\infty f(x - p) g(p) dp.
\]

Comportarea transformatei Fourier față de produsul de convoluție este dată de:
\[
  \kal{F}(f \ast g) = \sqrt{2\pi} \kal{F}(f) \cdot \kal{F}(g).
\]

Echivalent, acest rezultat se mai poate scrie prin inversare:
\[
  (f \ast g)(x) = \int_{-\infty}^\infty \widehat{f}(w) \cdot \widehat{g}(w) e^{iwx} dw.
\]

%%%%%%%%%%%%%%%%%%%%%%%%%%%%%%%%%%%%%%%%%%%%%%%%%%%%%%%%%%%%%%%%%%%%%% 

\section*{Exerciții}

Calculați transformatele Fourier ale următoarelor funcții:
\begin{enumerate}[(a)]
\item $ f(x)
  = \begin{cases}
    e^{2ix}, & -1 < x < 1 \\
    0, & \text{ în rest }
  \end{cases}
  $
\item $ f(x) =
  \begin{cases}
    1, & a < x < b \\
    0, & \text{ în rest }
  \end{cases}
  $
\item $ f(x) =
  \begin{cases}
    e^{kx}, & x < 0 \\
    0, & x > 0
  \end{cases}, k > 0
  $
\item $ f(x) =
  \begin{cases}
    e^x, & -a < x < a \\
    0, & \text{ în rest}
  \end{cases}
  $
\item $ f(x) = e^{-|x|}, x \in \RR $;
\item $ f(x) =
  \begin{cases}
    x, & 0 < x < a \\
    0, & \text{ în rest}
  \end{cases}
  $
\item $ f(x) =
  \begin{cases}
    xe^{-x}, & -1 < x < 0 \\
    0, \text{ în rest}
  \end{cases}
  $
\item $ f(x) =
  \begin{cases}
    |x|, & -1 < x < 1 \\
    0, & \text{ în rest}
  \end{cases}
  $
\item $ f(x) =
  \begin{cases}
    x, & -1 < x < 1 \\
    0, \text{ în rest}
  \end{cases}
  $
\end{enumerate}

%%%%%%%%%%%%%%%%%%%%%%%%%%%%%%%%%%%%%%%%%%%%%%%%%%%%%%%%%%%%%%%%%%%%%% 
\newpage

\section*{Tabele de transformate Fourier}

\begin{figure}[!htbp]
  \centering
  \begin{tabular}{c | c}
    $ f(x) $ & $ \widehat{f}_c(w) = \kal{F}_c(f) $ \\
    \hline\hline
    $ \begin{cases} 1, & 0 < x < a \\ 0, &\text{ în rest} \end{cases} $ & $ \sqrt{\frac{2}{\pi}} \frac{\sin aw}{w} $ \\
             & \\
    $ x^{a-1}, 0 < a < 1 $ & $ \sqrt{\frac{2}{\pi}} \frac{\Gamma(a)}{w^a} \cos \frac{a\pi}{2} $ \\
             & \\
    $ e^{-ax}, a > 0 $ & $ \sqrt{\frac{2}{\pi}} \cdot \frac{a}{a^2 + w^2} $ \\
             & \\
    $ e^{-x^2/2} $ & $ e^{-w^2/2} $ \\
             & \\
    $ e^{-ax^2}, a > 0 $ & $ \frac{1}{\sqrt{2a}} e^{-w^2/(4a)} $ \\
             & \\
    $ x^n e^{-ax}, a > 0 $ & $ \sqrt{\frac{2}{\pi}} \frac{n!}{(a^2 + w^2)^{n+1}} \mathrm{Re}(a + iw)^{n+1} $ \\
             & \\
    $ \begin{cases} \cos x, & 0 < x < a \\ 0, & \text{ în rest} \end{cases} $ & $ \frac{1}{\sqrt{2\pi}} \Big[ \frac{\sin a(1 - w)}{1 - w} + \frac{\sin a(1 + w)}{1 + w} \Big] $ \\
             & \\
    $ \cos (ax^2), a > 0 $ & $ \frac{1}{\sqrt{2a}} \cos \Big( \frac{w^2}{4a} - \frac{\pi}{4} \Big) $ \\
             & \\
    $ \sin (ax^2), a > 0 $ & $ \frac{1}{\sqrt{2a}} \cos \Big( \frac{w^2}{4a} + \frac{\pi}{4} \Big) $ \\
             & \\
    $ \frac{\sin ax}{x}, a > 0 $ & $ \sqrt{\frac{\pi}{2}} (1 - u(w - a)) $\footnotemark \\
             & \\
    $ \frac{e^{-x} \sin x}{x} $ & $ \frac{1}{\sqrt{2\pi}} \arctan \frac{2}{w^2} $ \\
  \end{tabular}
  \caption{\emph{Transformate Fourier cu cosinusuri (pentru funcții pare)}}
\end{figure}


\footnotetext{
  $ u(t - a) $ este funcția Heaviside (eng. \emph{unit step function}), definită prin:
  \[
    u(t - a) = \begin{cases}
      0, & t < a \\
      1, & t > a
    \end{cases}
  \]}

\newpage

\begin{figure}[!htbp]
  \centering
  \begin{tabular}{c|c}
    $ f(x) $ & $ \widehat{f}_s(x) = \kal{F}_s(f) $ \\
    \hline\hline
    $ \begin{cases} 1, & 0 < x < a \\ 0, & \text{ în rest} \end{cases} $ & $ \frac{\sqrt{2}}{\pi} \Big[ \frac{1 - \cos aw}{w} \Big] $ \\
             & \\
    $ \frac{1}{\sqrt{x}} $ & $ \frac{1}{\sqrt{w}} $ \\
             & \\
    $ x^{-3/2} $ & $ 2 \sqrt{w} $ \\
             & \\
    $  x^{a-1}, 0 < a < 1 $ & $ \sqrt{\frac{2}{\pi}} \frac{\Gamma(a)}{w^a} \sin\frac{a\pi}{2} $ \\
             & \\
    $ e^{-ax}, a > 0 $ & $ \sqrt{\frac{2}{\pi}} \cdot \frac{w}{a^2 + w^2} $ \\
             & \\
    $ \frac{e^{-ax}}{x}, a > 0 $ & $ \sqrt{\frac{2}{\pi}} \arctan \frac{w}{a} $ \\
             & \\
    $ x^n e^{-ax}, a > 0 $ & $ \sqrt{\frac{2}{\pi}} \frac{n!}{(a^2 + w^2)^{n+1}} \mathrm{Im}(a + iw)^{n + 1} $ \\
             & \\
    $ xe^{-x^2/2} $ & $ we^{-w^2/2} $ \\
             & \\
    $ xe^{-ax^2}, a > 0 $ & $ \frac{w}{(2a)^{3/2}} e^{-w^2/(4a)} $ \\
             & \\
    $ \begin{cases} \sin x, & 0 < x < a \\ 0, &\text{ în rest} \end{cases} $ & $ \frac{1}{\sqrt{2\pi}} \Big[ \frac{\sin a(1 - w)}{1 - w} - \frac{\sin a(1 + w)}{1 + w} \Big] $ \\
             & \\
    $ \frac{\cos ax}{x}, a > 0 $ & $ \sqrt{\frac{\pi}{2}} u(w - a) $ \\
             & \\
    $ \arctan{2a}{x}, a > 0 $ & $ \sqrt{2\pi} \frac{\sin aw}{w} e^{-aw} $ 
  \end{tabular}
  \caption{\emph{Transformate Fourier cu sinusuri (pentru funcții impare)}}
\end{figure}

\newpage

\begin{figure}[!htbp]
  \centering
  \begin{tabular}{c|c}
    $ f(x) $ & $ \widehat{f}(w) = \kal{F}(f) $ \\
    \hline\hline
    $ \begin{cases} 1, & -b < x < b \\ 0, &\text{ în rest} \end{cases} $ & $ \sqrt{\frac{2}{\pi}} \frac{\sin bw}{w} $ \\
             & \\
    $ \begin{cases} 1, & b < x < c \\ 0, &\text{ în rest} \end{cases} $ & $ \frac{e^{-ibw} - e^{-icw}}{iw \sqrt{2\pi}} $ \\
             & \\
    $ \frac{1}{x^2 + a^2}, a > 0 $ & $ \sqrt{\frac{\pi}{2}} \frac{e^{-a|w|}}{a} $ \\
             & \\
    $ \begin{cases} e^{-ax}, & x > 0 \\ 0, &\text{ în rest} \end{cases}, a > 0 $ & $ \frac{1}{\sqrt{2\pi}}{a + iw} $ \\
    $ \begin{cases} e^{ax}, & b < x < c \\ 0, &\text{ în rest} \end{cases} $ & $ \frac{e^{(a-iw)c} - e^{(a - iw)b}}{\sqrt{2\pi}(a - iw)} $ \\
             & \\
    $ \begin{cases} e^{iax}, & -b < x < b \\ 0, &\text{ în rest} \end{cases} $ & $ \sqrt{\frac{2}{\pi}} \frac{\sin b(w - a)}{w - a} $ \\
             & \\
    $ \begin{cases} e^{iax}, & b < x < c \\ 0, &\text{ în rest} \end{cases} $ & $ \frac{i}{\sqrt{2\pi}} \frac{e^{ib(a - w)} - e^{ic(a - w)}}{a - w} $ \\
             & \\
    $ e^{-ax^2}, a > 0 $ & $ \frac{1}{\sqrt{2a}} e^{-w^2/(4a)} $ \\
             & \\
    $ \frac{\sin ax}{x}, a > 0 $ & $ \begin{cases} \sqrt{\frac{\pi}{2}}, & |w| < a \\ 0, & |w| > a \end{cases} $
  \end{tabular}
  \caption{\emph{Transformate Fourier generale (pentru funcții arbitrare)}}
\end{figure}


\end{document}
